% ============================================================
% SECTION 4.2: WHAT THE ACCOUNTING COSTS. Assembled 2026-07-26 (late).
%
% RELOCATION, not rewriting. Every table and every number here is moved verbatim
% from sections/_mined/03_protocol_v1_full.tex, which held them while the
% protocol section was still carrying its own evidence. Only the connective
% prose is new, and its job is to make one argument out of what used to be six
% separate control demonstrations.
%
% Provenance of each figure is unchanged and was verified against CSVs earlier:
%   uniform table      exp_G/results.csv (arm=artifact1, 3 seeds, means)
%   63-cell table      exp_P/results.csv (dQ_naive, dQ_honest_vs_mean)
%   granularity table  exp_P/recovery.csv (com-DBLP, louvain)
%   compute inversion  exp_P/results.csv + bestof.csv; exp_C/results_summary.csv
%   permuted weights   exp_V/results.csv + null_arm.csv
%   retention audit    exp_J_figures/clip_fractions.csv; exp_C; exp_B
% STYLE: no em-dashes, no AI-characteristic phrasing (see STYLE.md).
% ============================================================

\subsection{What the accounting costs}\label{sec:experiments:accounting}

Each control in Section~\ref{sec:protocol} exists because a comparison can favour
the pipeline without any improvement having occurred. This subsection measures
how much each one is worth. The measurements are of the evaluation and not of the
sparsifiers, and they set the size of the correction that the rest of the paper
applies.

\paragraph{Scoring on the sparsified graph}
Table~\ref{tab:uniform} takes uniform random edge deletion, whose selection rule
uses no information whatever, and scores it both ways on three networks at two
retentions. Read off the sparsified graph it improves modularity in all six
cells. The same partitions scored on the original graph are worse than the
baseline in all six. The sign of the conclusion is opposite in every cell, and
the gap widens as retention falls, which is what
Section~\ref{sec:protocol:transfer} predicts: the sparser the graph, the higher
the modularity an optimizer can reach on it.

\begin{table}[t]
\centering
\caption{Uniform random edge deletion under the two accountings, Leiden, mean of
three seeds. $\Delta Q_{\mathrm{self}}$ is positive in every cell and
$\Delta Q_{\mathrm{transfer}}$ is negative in every cell. Uniform deletion uses
no structural information, so the inflation comes from the objective rather than
from any selection rule.}
\label{tab:uniform}
\small
\setlength{\tabcolsep}{4pt}
\begin{tabular}{@{}lccrr@{}}
\toprule
Network & $Q_{\mathrm{base}}$ & $\rho$ & $\Delta Q_{\mathrm{self}}$ &
$\Delta Q_{\mathrm{transfer}}$\\
\midrule
ca-CondMat  & 0.729 & 0.800 & $+0.0111$ & $-0.0164$\\
ca-CondMat  & 0.729 & 0.900 & $+0.0053$ & $-0.0064$\\
email-Enron & 0.611 & 0.800 & $+0.0052$ & $-0.0162$\\
email-Enron & 0.611 & 0.900 & $+0.0011$ & $-0.0086$\\
com-DBLP    & 0.826 & 0.800 & $+0.0078$ & $-0.0247$\\
com-DBLP    & 0.826 & 0.900 & $+0.0040$ & $-0.0098$\\
\bottomrule
\end{tabular}
\end{table}

The effect is not particular to that sparsifier or that detector.
Table~\ref{tab:accounting} covers 63 configurations spanning three detection
algorithms, seven networks and three sparsifier settings. Modularity read off the
sparsified graph is positive in 60 of them and reaches $+0.537$; the same
partitions scored on the original graph are positive in 15. Between the two
accountings the sign of the conclusion differs in 45 of the 63 configurations.
Eleven cells exceed a compute-matched baseline, all of them under label
propagation, and they are not a subset of the 15: three sit below the plain
baseline mean while beating a matched control that had drawn a single restart,
which is the failure the compute control addresses.

\begin{table}[t]
\centering
\caption{The two accountings across 63 configurations, seven networks by three
sparsifier settings under each algorithm. Both quantities are measured against
the mean of the plain restarts, the weakest of the baselines in
Section~\ref{sec:protocol:compute}; counts against the stricter best-of-$N$
baseline appear in Section~\ref{sec:experiments:negative}. A flip is a cell where
$\Delta Q_{\mathrm{self}} > 0$ and $\Delta Q_{\mathrm{transfer}} < 0$, so that
the two accountings disagree on the sign of the conclusion.}
\label{tab:accounting}
\small
\setlength{\tabcolsep}{4pt}
\begin{tabular}{@{}lcccc@{}}
\toprule
 & Infomap & Louvain & Lab.\ prop. & Total\\
\midrule
Cells                              & 21 & 21 & 21 & 63\\
$\Delta Q_{\mathrm{self}} > 0$     & 20 & 20 & 20 & 60\\
$\Delta Q_{\mathrm{transfer}} > 0$ & 6  & 0  & 9  & 15\\
Flips                              & 14 & 20 & 11 & 45\\
max $\Delta Q_{\mathrm{self}}$ & $+0.312$ & $+0.260$ & $+0.537$ & $+0.537$\\
\bottomrule
\end{tabular}
\end{table}

\paragraph{Granularity}
Table~\ref{tab:granularity} gives the anatomy of a recovery gain on com-DBLP
under Louvain. L-Spar at realized retention $0.503$ raises average best-match
$F_1$ from $0.102$ to $0.193$, an apparent gain of $+0.091$, while taking the
number of communities of at least three members from 220 to 10{,}947.
Partitioning the original graph to the same count reaches $0.300$. The finer
partition is worth more than twice what the sparsification appeared to be worth
and requires no sparsifier. The same pattern holds for the degree-based arm on
the same network, where the apparent change is negative and the matched control
is positive. Leiden reproduces the sparsified figure to four decimal places and
its own control, at $\gamma = 576$ and 10{,}812 communities, reaches $0.297$, so
the effect belongs to the modularity objective's default granularity and not to
one optimizer's search.

\begin{table}[t]
\centering
\caption{Granularity control on com-DBLP under Louvain, average best-match $F_1$
over reference communities of size at least three. $k_{\geq 3}$ counts
communities of that size; $\gamma$ is the resolution used on the \emph{original}
graph. Each sparsified arm is followed by its control, a partition of the
original graph at the same $k_{\geq 3}$. The floor is a size-matched random
partition at the arm's own $k_{\geq 3}$. Baseline and sparsified rows are means
over three to six seeds; controls are single runs.}
\label{tab:granularity}
\small
\setlength{\tabcolsep}{4pt}
\begin{tabular}{@{}lrrcc@{}}
\toprule
Condition & $k_{\geq 3}$ & $\gamma$ & $F_1$ & floor\\
\midrule
Original graph               & 220    & 1   & $0.102 \pm 0.005$ & 0.019\\
DSpar, $\rho = 0.500$        & 1{,}684  & 1   & $0.098 \pm 0.002$ & 0.049\\
\quad matched control        & 1{,}657  & 24  & $0.125$           & \\
L-Spar, $\rho = 0.503$       & 10{,}947 & 1   & $0.193 \pm 0.000$ & 0.090\\
\quad matched control        & 11{,}159 & 640 & $0.300$           & \\
\bottomrule
\end{tabular}
\end{table}

Where the counts cannot be brought together, we draw no recovery conclusion. The
comparability rule is $|\Delta k|/k < 25\%$, and the cells it excludes are
reported as excluded rather than compared.

\paragraph{Chance}
The floor column of Table~\ref{tab:granularity} carries the second artifact. A
random partition with the same community size distribution scores $0.019$ against
the reference communities at 220 communities and $0.090$ at 10{,}947, so a
comparison that reads $+0.091$ from a fragmenting sparsifier is reading a floor
that rose by $0.071$ underneath it. The same measurement on com-Amazon gives
$0.009$ at 231 communities and $0.057$ at 8{,}651, a $6.6\times$ rise from
fragmentation alone.

\paragraph{The compute budget}
The registered compute-matched control buys less than it appears to. In a sweep
of 48 cells over six networks, four retention targets and two samplers under
Leiden, it bought exactly two restarts in every cell; across the 63
configurations of Table~\ref{tab:accounting} it bought between 1 and 16, and 1 in
13 of the 18 cells on the two largest networks. Label propagation shows the
consequence. Under the compute-matched control it is the strongest of the three
algorithms in that table, with a mean change of $+0.009$; against the best of 20
plain restarts of the same baseline, 10 on the two largest networks, it is the
weakest, at $-0.106$. On email-Enron its two apparent gains of $+0.146$ and
$+0.161$ become $-0.068$ and $-0.053$. The compute-matched baseline drew four and
five restarts in those two cells and still never sampled the algorithm's good
mode, whose modularity is $0.523$ against the matched control's $0.308$.

\paragraph{Weighting credited to similarity}
On email-Enron, weighting each edge by $1 + J(u,v)$, with $J$ the Jaccard
similarity of the endpoint neighbourhoods, raises modularity measured on the
original graph to $0.6153 \pm 0.0034$ against a baseline of $0.6071$ and a
compute-matched best of $0.6051$, a gain of $+0.0102$ at $3.7$ times the pooled
standard deviation. Permuting those same weights across edges, which preserves
their distribution and destroys their placement, gives $0.6168 \pm 0.0020$, a
gain of $+0.0117$ at $6.0$ pooled standard deviations. The null reproduces the
gain and exceeds it. On a degree-preserving rewiring of ca-CondMat, where the
similarity signal is absent by construction, weighting still produces $+0.0023$
and clears the same significance test. What the weights contribute is
heterogeneity in the optimizer's search and not information about communities.

\paragraph{Nominal retention}
The parameter a sampler is given does not determine what it deletes. The scheme
in common use clips per-edge probabilities at one, and across six networks at
four targets between 12\% and 36\% of edges reach that clip, so the expected
number retained falls below the nominal $\alpha m$ and saturates: on
email-Eu-core the expected retention at $\alpha = 1.0$ is $0.665$, and on
wiki-Vote at $\alpha = 0.95$ it is $0.447$. Three samplers that describe
themselves by the same $\alpha$ realize different fractions. Sampling with
replacement gives $0.33$ to $0.52$ across seventeen networks at a nominal $0.8$,
the clipped no-replacement scheme gives $0.37$ to $0.68$ across those six
networks and four targets, and solving the scale factor by bisection on the same
cells gives $0.70$ to $0.95$. A comparison at equal nominal parameters is a
comparison at unequal deletion.
